\documentclass[11pt]{article}
\usepackage[margin=1in]{geometry}
\usepackage{amsmath,booktabs,array,graphicx,pdflscape,hyperref}
\hypersetup{hidelinks}
\setlength{\emergencystretch}{3em}
\providecommand{\Description}[1]{}
\ifdefined\pdfinfoomitdate\pdfinfoomitdate=1\fi
\ifdefined\pdfsuppressptexinfo\pdfsuppressptexinfo=-1\fi
\ifdefined\pdftrailerid\pdftrailerid{}\fi
\renewcommand{\thetable}{S\arabic{table}}
\let\originaltable\table
\let\endoriginaltable\endtable
\renewenvironment{table}[1][]{\originaltable[!htbp]}{\endoriginaltable}
\title{Technical Supplement\\Emergency Review and Constitutional Court Design by Simulation}
\author{Anonymous Author(s)}
\date{}
\begin{document}
\maketitle

This supplement preserves detailed diagnostic tables and representative model weights accompanying the anonymous manuscript. The quantities are synthetic model outputs or coding assumptions, not estimates of causal reform effects. The replication materials contain the full source, frozen inputs, and report generators.

Table~\ref{tab:normative-scores} reports alternative reading aids for different constitutional priorities: rights priority, emergency restraint, democratic responsiveness, legal stability, low conflict, and administrative feasibility. The table makes visible that any single score embeds a normative weighting choice.
% Auto-generated by paper/scripts/generate_figures.py
\begin{table}[p]
\centering
\caption{Multi-objective reading aids under different normative priorities}
\label{tab:normative-scores}
\Description{Table reporting alternative synthetic score families for rights priority, emergency restraint, democratic responsiveness, legal stability, low conflict, and administrative feasibility.}
\scriptsize
\setlength{\tabcolsep}{3pt}
\resizebox{\textwidth}{!}{%
\begin{tabular}{>{\raggedright\arraybackslash}p{1.45in}rrrrrr}
\toprule
Scenario & Rights priority & Emergency restraint & Democratic response & Legal stability & Low conflict & Admin feasibility \\
\midrule
Current-like & 0.475 & 0.437 & 0.591 & 0.609 & 0.700 & 0.668 \\
18-year terms & 0.455 & 0.627 & 0.629 & 0.620 & 0.730 & 0.691 \\
Written reasons & 0.450 & 0.620 & 0.629 & 0.625 & 0.726 & 0.689 \\
Auto merits & 0.506 & 0.685 & 0.640 & 0.598 & 0.734 & 0.691 \\
No merits gap & 0.495 & 0.690 & 0.646 & 0.605 & 0.738 & 0.696 \\
Recusal enforce. & 0.457 & 0.627 & 0.635 & 0.619 & 0.733 & 0.679 \\
Judicial electorate & 0.456 & 0.628 & 0.636 & 0.620 & 0.734 & 0.674 \\
Random panels & 0.452 & 0.627 & 0.639 & 0.622 & 0.731 & 0.661 \\
Integrity package & 0.500 & 0.685 & 0.647 & 0.603 & 0.738 & 0.684 \\
Constitutional remand & 0.436 & 0.629 & 0.641 & 0.651 & 0.742 & 0.672 \\
Public-interest filter & 0.465 & 0.629 & 0.634 & 0.624 & 0.734 & 0.681 \\
\bottomrule
\end{tabular}
}%
\end{table}


The directional score is intentionally secondary:
\[
\begin{aligned}
\mathrm{DirectionalScore}_s = \mathrm{mean}(&\mathrm{StabilityRights}_s,
\mathrm{LegitimacyControl}_s,\mathrm{ClaimantSuccess}_s,\\
&\mathrm{PrecedentDurability}_s,\mathrm{LowerCourtCompliance}_s,
\mathrm{EliteAcceptance}_s,\\
1-\mathrm{AdministrativeCost}_s).
\end{aligned}
\]
The stability-rights component averages legal stability, precedent durability, lower-court compliance, enforcement capacity, lower-court-resistance restraint, rights protection, reversal restraint, and conflict restraint. The legitimacy-control component averages legitimacy, democratic responsiveness, independence-accountability balance, partisan-alignment restraint, emergency-irregularity restraint, emergency-legitimacy-risk restraint, emergency-downstream restraint, government-noncompliance restraint, recusal-pressure restraint, forum-shopping restraint, emergency-opportunism restraint, and strategic-pressure restraint. Because those choices are normative, the paper uses the directional score only as a reading aid. Score differences below .010 are treated as indistinguishable in the manuscript's substantive interpretation.

\clearpage
\begin{landscape}
\setlength{\textwidth}{\linewidth}
Table~\ref{tab:pipeline-diagnostics} is the audit table for the concern that formal court design might overweight the model. It exposes litigation-pipeline and downstream-enforcement quantities directly: certiorari admission, court-requested response, CVSG request, bar capital, claim strength, vehicle quality, genuine lower-court split rate, lower-court split depth, lower-court resistance, forum shopping, pre-review settlement, strategic plaintiff selection, repeat-player advantage, enforcement capacity, emergency opportunism, recusal incentive pressure, government noncompliance, emergency downstream effect, and precedent durability.
% Auto-generated by paper/scripts/generate_figures.py
\begin{table}[hbt!]
\centering
\caption{Litigation-pipeline and downstream-enforcement diagnostics}
\label{tab:pipeline-diagnostics}
\Description{Table reporting certiorari admission, lower-court split depth, lower-court resistance, forum shopping, settlement, strategic plaintiff selection, repeat-player advantage, enforcement capacity, emergency opportunism, recusal incentive pressure, government noncompliance, emergency downstream effect, and precedent durability for selected designs.}
\scriptsize
\setlength{\tabcolsep}{2pt}
\resizebox{\textwidth}{!}{%
\begin{tabular}{>{\raggedright\arraybackslash}p{1.38in}rrrrrrrrrrrrr}
\toprule
Scenario & Cert admit & Split & Resistance & Forum & Settlement & Plaintiff sel. & Repeat player & Enforcement & Emerg. opp. & Recusal press. & Gov. noncomp. & Emerg. down. & Prec. dur. \\
\midrule
Current-like & 0.584 & 0.589 & 0.327 & 0.483 & 0.005 & 0.487 & 0.261 & 0.547 & 0.541 & 0.470 & 0.188 & 0.196 & 0.679 \\
18-year terms & 0.582 & 0.589 & 0.316 & 0.444 & 0.005 & 0.487 & 0.261 & 0.559 & 0.384 & 0.477 & 0.170 & 0.122 & 0.722 \\
Written emergency reasons & 0.581 & 0.589 & 0.316 & 0.444 & 0.006 & 0.487 & 0.261 & 0.559 & 0.357 & 0.477 & 0.171 & 0.133 & 0.739 \\
Automatic merits follow-up & 0.583 & 0.589 & 0.312 & 0.443 & 0.005 & 0.487 & 0.261 & 0.564 & 0.335 & 0.477 & 0.167 & 0.118 & 0.703 \\
No emergency merits gap & 0.600 & 0.589 & 0.312 & 0.444 & 0.006 & 0.487 & 0.261 & 0.565 & 0.316 & 0.492 & 0.168 & 0.106 & 0.736 \\
Strong recusal enforcement & 0.599 & 0.589 & 0.316 & 0.444 & 0.005 & 0.487 & 0.261 & 0.560 & 0.384 & 0.498 & 0.176 & 0.125 & 0.727 \\
Random merits panels & 0.600 & 0.589 & 0.316 & 0.414 & 0.006 & 0.487 & 0.261 & 0.560 & 0.384 & 0.620 & 0.172 & 0.126 & 0.737 \\
Emergency integrity package & 0.598 & 0.589 & 0.312 & 0.444 & 0.006 & 0.487 & 0.261 & 0.565 & 0.336 & 0.497 & 0.169 & 0.121 & 0.734 \\
Jurisdiction-stripping constraints & 0.594 & 0.589 & 0.316 & 0.394 & 0.005 & 0.487 & 0.261 & 0.590 & 0.384 & 0.477 & 0.183 & 0.124 & 0.729 \\
Legislative override window & 0.595 & 0.589 & 0.317 & 0.444 & 0.006 & 0.487 & 0.261 & 0.561 & 0.384 & 0.478 & 0.182 & 0.124 & 0.728 \\
Constitutional remand & 0.625 & 0.589 & 0.282 & 0.405 & 0.006 & 0.487 & 0.261 & 0.580 & 0.385 & 0.483 & 0.178 & 0.128 & 0.779 \\
Public-interest filter & 0.634 & 0.589 & 0.316 & 0.361 & 0.006 & 0.487 & 0.261 & 0.562 & 0.385 & 0.482 & 0.185 & 0.131 & 0.731 \\
Remand + override window & 0.627 & 0.589 & 0.282 & 0.405 & 0.006 & 0.487 & 0.261 & 0.580 & 0.358 & 0.482 & 0.181 & 0.139 & 0.789 \\
Panel + jurisdiction safeguards & 0.602 & 0.589 & 0.317 & 0.364 & 0.006 & 0.487 & 0.261 & 0.588 & 0.371 & 0.657 & 0.184 & 0.140 & 0.758 \\
\bottomrule
\end{tabular}
}%
\end{table}

\end{landscape}

\clearpage
Table~\ref{tab:mechanical-emergent} records the interpretive rule used in the results section: index movement is most informative when it travels with downstream rights, compliance, implementation, conflict, or cost measures.
% Auto-generated by paper/scripts/generate_figures.py
\begin{table}[hbt!]
\centering
\caption{Mechanical implications versus diagnostic findings}
\label{tab:mechanical-emergent}
\Description{Table distinguishing model implications that are partly coded into metric definitions from findings that depend on downstream simulated tradeoffs.}
\footnotesize
\setlength{\tabcolsep}{4pt}
\resizebox{\textwidth}{!}{%
\begin{tabular}{>{\raggedright\arraybackslash}p{1.7in}>{\raggedright\arraybackslash}p{1.0in}>{\raggedright\arraybackslash}p{2.25in}>{\raggedright\arraybackslash}p{2.1in}}
\toprule
Claim & Mechanical status & What would be diagnostic & Manuscript use \\
\midrule
Written emergency reasons reduce emergency irregularity & Partly coded & Diagnostic only if downstream disruption, compliance, or rights outputs also move & Use emergency profile and mechanism contrasts; do not report as independent causal discovery \\
Automatic merits follow-up reduces merits-displacing emergency relief & Partly coded & Emergent only through workload, lower-court compliance, and public-legitimacy proxy changes & Read with emergency walkthrough and mechanism table \\
No emergency relief without merits review produces the lowest irregularity & Strongly coded & Tradeoff is whether rights claimant success, settlement, or downstream compliance worsens & Treat as design implication plus tradeoff test \\
Judicial-electorate selection affects appointment-pressure diagnostics more than emergency outputs & Not mechanically forced by the emergency metric & Emergent if legitimacy/alignment changes without comparable emergency movement & Use as comparison mechanism, not main emergency-reform result \\
Randomized panels or remand rules have mixed profiles & Weakly coded & Emergent through access, correction, compliance, and administrative-cost channels & Use as evidence that non-emergency mechanisms move costs elsewhere \\
\bottomrule
\end{tabular}
}%
\end{table}


\clearpage
Table~\ref{tab:model-weights} records representative weights used by the simulator. These are source-informed and design-prior coding choices, not fitted coefficients.
\begin{table}[p]
\centering
\caption{Representative model weights and decision rules}
\label{tab:model-weights}
\Description{Table summarizing representative admission, voting, emergency, conflict, legitimacy, and override weights used by the simulator.}
\scriptsize
\renewcommand{\arraystretch}{0.9}
\begin{tabular}{>{\raggedright\arraybackslash}p{1.10in}>{\raggedright\arraybackslash}p{2.65in}>{\raggedright\arraybackslash}p{2.05in}}
\toprule
Component & Representative rule & Interpretation \\
\midrule
Admission & $0.18 + .22$ cert pressure $+ .12$ lower-court conflict $+ .10$ lower-court error $+ .13$ split depth $+ .06$ genuine split $+ .15$ SG signal $+ .10$ amicus signal $+ .08$ split maturity $+ .07$ relist signal $+ .06$ bar capital $+ .10$ claim strength $+ .09$ vehicle quality $+ .05$ specialist counsel $+ .05$ repeat-player signal $+ .05$ strategic-plaintiff signal $+ .02$ forum-shopping signal $+ .08$ conditional reversal $-.06$ settlement pressure $-.16$ vehicle defect, plus path adjustment, CFR, and CVSG increments & Admission is an institutional screen, not a fixed docket quota. \\
Voting & Legal concern receives weight $.42$; rights concern uses rights burden, justice rights sensitivity, and weight $.44$; claim strength, vehicle quality, and bar capital can modestly increase merits credibility; democratic mandate, legislative quality, institutionalism, emergency deference, repeat-player pressure, forum-shopping signal, precedent stability, and accountability pull reduce or increase invalidation propensity. & Merits votes mix legal, rights, strategic, institutional, and emergency-deference considerations. \\
Emergency irregularity & Open emergency rule base $.32$, no-merits penalty $.22$, shadow stay $.24$, reasoned order $-.08$, merits acceleration $-.07$, partisan salience $.12$, ambiguity $.08$, emergency opportunism $.10$. & The legacy \texttt{shadowDocketAbuse} field is retained in CSV outputs for compatibility; the paper-facing index is emergency irregularity. \\
Process legitimacy & Legitimacy base $.48$, public attention $.16$, legislative quality $.10$, merits reason-giving $.18$, recusal discipline up to $.16$, council warning $.06$, emergency-process boost or shadow-stay penalty, emergency irregularity $-.24$, partisan alignment $-.20$; the process-legitimacy proxy then subtracts emergency legitimacy risk $.18$, partisan alignment $.12$, appointment pressure $.08$, and court-curbing $.10$. & The index rewards visible process and penalizes opaque or partisan-looking intervention; it is not direct public-opinion measurement. \\
Compliance & Lower-court compliance rises with enforcement capacity and merits clarity, and falls with lower-court conflict, error risk, split depth, ideological drift, lower-court resistance, state conflict, invalidation shock, and government noncompliance. & Enforcement and implementation are modeled separately from formal legal validity. \\
Conflict & Conflict potential $.38$, executive defiance $.18$, lower-court conflict $.08$, lower-court resistance $.10$, forum shopping $.04$, conflict load $.12$, legislative defiance $.10$, government noncompliance, emergency downstream effect, invalidating high-mandate laws $.14$, emergency relief $.12$, cross-court disagreement $.16$, override attempt $.06$. & Constitutional conflict is cumulative and rises when legal disagreement becomes interbranch disagreement. \\
Override & Attempt pressure combines override pressure $.34$, democratic mandate $.30$, public attention $.18$, conflict load $.12$, legislative defiance $.12$, override adaptation $.10$, and rights-burden penalty $-.24$. & Accountability devices affect implementation only after invalidation and can generate rights carveouts or repeated override loops. \\
\bottomrule
\end{tabular}
\end{table}


\clearpage
\end{document}
